\documentclass[a4paper,11pt]{article}
\usepackage[a4paper,margin=2cm]{geometry}
\usepackage[T1]{fontenc}
\usepackage[utf8]{inputenc}
\usepackage[ngerman,english]{babel}
\usepackage{lmodern}
\usepackage{booktabs,longtable,array}
\usepackage[hidelinks]{hyperref}
\title{Adjektivellipse\\\large Omission of a recoverable noun after an adjective}
\date{}
\begin{document}
\maketitle
\tableofcontents

\section{Definition und Grundprinzip}
Eine \textbf{Adjektivellipse} ist eine elliptische Nominalgruppe: Das Nomen wird ausgelassen, weil es aus dem sprachlichen oder situativen Kontext erschließbar ist. Das Adjektiv und gegebenenfalls ein Artikelwort bleiben stehen und übernehmen die Funktion der Nominalgruppe.

\textbf{English:} An adjective ellipsis omits a recoverable noun while keeping the adjective (and possibly a determiner).

\textit{Ich habe zwei Hemden. Das blaue [Hemd] gefällt mir besser.}\\
\textit{I have two shirts. I like the blue one better.}

\textbf{Wichtig:} Nicht jede Form ohne sichtbares Nomen ist automatisch eine Adjektivellipse. Insbesondere ist zwischen Ellipse, substantiviertem Adjektiv und prädikativem Adjektiv zu unterscheiden.

\section{Wann ist die Auslassung möglich?}
\begin{itemize}
\item Das fehlende Nomen ist durch eine frühere Erwähnung klar: \textit{Welchen Mantel willst du? -- Den schwarzen [Mantel].}
\item Die Situation liefert das Nomen: \textit{Ich nehme den roten [Pullover].} (Beim Einkauf mit sichtbaren Pullovern.)
\item Das ausgelassene Nomen muss möglichst eindeutig rekonstruierbar sein. Bei Mehrdeutigkeit wird es besser wiederholt.
\item Ellipsen sind häufig bei Auswahl, Vergleichen und Wiederholungsvermeidung.
\end{itemize}

\section{Kasus, Genus und Numerus}
Das ausgelassene Nomen bestimmt weiterhin \textbf{Genus und Numerus}; die syntaktische Funktion bzw. Präposition bestimmt den \textbf{Kasus}. Das verbliebene Adjektiv wird weiterhin dekliniert.

\begin{longtable}{p{2.5cm}p{5.2cm}p{5.2cm}}
\toprule
Kasus & Vollständige Nominalgruppe & Adjektivellipse\\\midrule
Nominativ & Der rote Mantel ist teuer. & Der rote ist teuer.\\
Akkusativ & Ich kaufe den roten Mantel. & Ich kaufe den roten.\\
Dativ & Ich helfe dem jungen Mann. & Ich helfe dem jungen. (kontextabhängig; häufig unnatürlich)\\
Genitiv & Die Farbe des roten Mantels gefällt mir. & Die Farbe des roten gefällt mir. (nur mit eindeutigem Bezug)\\
\bottomrule
\end{longtable}
\textbf{Beachte:} Nicht jede grammatisch mögliche Ellipse ist idiomatisch. Besonders bei Personen oder schwer erschließbaren Referenten ist das volle Nomen oft besser.

\section{Deklination: starke, schwache und gemischte Formen}
Die üblichen Deklinationsregeln für attributive Adjektive gelten grundsätzlich auch beim elliptisch ausgelassenen Nomen:
\begin{itemize}
\item \textbf{Schwach nach bestimmtem Artikel:} \textit{der rote [Mantel], den roten [Mantel], das blaue [Hemd], die grünen [Schuhe]}.
\item \textbf{Gemischt nach ein-/Possessivwörtern:} \textit{ein roter [Mantel], ein blaues [Hemd], mein neuer [Computer], meine alten [Schuhe]}.
\item \textbf{Stark ohne Artikel:} \textit{Roter [Wein] ist noch da; weißer [Wein] ist aus.} Solche Fälle können auch als feste Verwendungen oder Substantivierungen analysiert werden; im normalen Sprachgebrauch ist ein explizites Bezugsnomen oder ein geeigneter Determinierer oft klarer.
\end{itemize}
Das Artikelwort allein ist nicht die gesamte Erklärung: Die Endung hängt zusätzlich von Kasus, Genus, Numerus und Deklinationstyp ab.

\section{Ein-, kein- und Possessivwörter: wichtige Besonderheit}
Wenn ein Nomen ausgelassen wird, können \textbf{ein/kein/sein/mein} selbständig als pronominale Formen erscheinen. Dadurch können sich die Formen von attributiven Artikelwörtern unterscheiden:
\begin{itemize}
\item \textit{Hast du einen Stift? -- Ja, ich habe \textbf{einen}.} (Indefinitpronomen, kein verbliebenes Adjektiv)
\item \textit{Hast du einen roten Stift? -- Ja, ich habe \textbf{einen roten}.} (Adjektivellipse)
\item \textit{Ich brauche einen neuen Stift. -- Ich brauche \textbf{einen neuen}.}
\item \textit{Ist das dein Buch? -- Nein, \textbf{meins}.} (selbständiges Possessivpronomen; keine Adjektivellipse)
\end{itemize}
Insbesondere \textit{ein roter} (vollständige Gruppe mit Nomen) und \textit{einen roten} (elliptische Gruppe im Akkusativ) dürfen nicht verwechselt werden.

\section{Abgrenzung zu ähnlichen Konstruktionen}
\subsection{Substantivierte Adjektive}
Substantivierte Adjektive funktionieren als Nomen und werden großgeschrieben: \textit{der Alte, die Neue, etwas Schönes, nichts Besonderes}. Sie müssen nicht auf ein zuvor genanntes, rekonstruierbares Nomen zurückgehen.

Bei \textbf{Adjektivellipse} bleibt das Adjektiv gewöhnlich \textbf{kleingeschrieben}, sofern es sich auf ein ausgelassenes konkretes Nomen bezieht: \textit{Welches Fahrrad? -- Das rote [Fahrrad].}

\textbf{Kontrast:}\\
\textit{Die neue [Jacke] ist teuer.} (Ellipse, Bezug: Jacke)\\
\textit{Die Neue arbeitet hier.} (Substantivierung, eine neue Mitarbeiterin)\\
\textit{Ich möchte das Gute bewahren.} (Substantivierung, kein konkretes ausgelassenes Nomen)

\subsection{Prädikative und adverbiale Adjektive}
\textit{Das Auto ist rot.} (prädikativ; keine Deklinationsendung, keine Adjektivellipse)\\
\textit{Er fährt schnell.} (adverbial; keine Deklinationsendung, keine Adjektivellipse)\\
\textit{Das rote [Auto] fährt schnell.} (attributives Adjektiv mit ausgelassenem Nomen)

\subsection{Ellipsen anderer Wortarten}
\textit{Ich nehme dieses.} (Demonstrativpronomen, kein erhaltenes Adjektiv)\\
\textit{Ich nehme das rote.} (Adjektivellipse)\\
\textit{Ich nehme das Rote.} (substantivierte Farbbezeichnung; abhängig von gemeinter Bedeutung)

\section{Typische Beispielsätze mit Erklärung}
\begin{enumerate}
\item \textit{Welche Tasche gefällt dir? -- Die schwarze [Tasche].}\\
\textbf{English:} Which bag do you like? -- The black one. (Nom. fem.)
\item \textit{Welchen Pullover kaufst du? -- Den blauen [Pullover].}\\
\textbf{English:} Which sweater are you buying? -- The blue one. (Acc. masc.)
\item \textit{Welches Auto gehört dir? -- Das kleine [Auto].}\\
\textbf{English:} Which car is yours? -- The small one. (Nom. neut.)
\item \textit{Mit welcher Kamera fotografierst du? -- Mit der neuen [Kamera].}\\
\textbf{English:} Which camera are you using? -- With the new one. (Dat. fem.)
\item \textit{Welche Bücher brauchst du? -- Die deutschen [Bücher].}\\
\textbf{English:} Which books do you need? -- The German ones. (Acc. plural.)
\item \textit{Ich nehme lieber den großen [Koffer] als den kleinen [Koffer].}\\
\textbf{English:} I prefer the large suitcase to the small one.
\end{enumerate}

\section{Häufige Fehler}
\begin{enumerate}
\item \textbf{Falscher Kasus:} \textit{Ich kaufe der rote.} $\rightarrow$ \textit{Ich kaufe den roten [Mantel].}
\item \textbf{Falsche Endung:} \textit{Das rot gefällt mir.} $\rightarrow$ \textit{Das rote [Fahrrad] gefällt mir.}
\item \textbf{Unnötige Großschreibung:} \textit{die Schwarze [Tasche]} $\rightarrow$ \textit{die schwarze [Tasche]}.
\item \textbf{Unklarer Bezug:} \textit{Ich nehme den grünen.} ist nur klar, wenn das Bezugsnomen bekannt ist.
\item \textbf{Verwechslung mit Prädikativ:} \textit{Der Pullover ist blau} ist keine Adjektivellipse.
\end{enumerate}

\section{Entscheidungshilfe}
\begin{enumerate}
\item Kann man ein \textbf{konkretes Nomen} aus dem Kontext ergänzen? Falls ja, liegt möglicherweise eine Adjektivellipse vor.
\item Steht das Adjektiv \textbf{prädikativ oder adverbial}? Dann liegt keine Adjektivellipse vor.
\item Ist es als eigenständiges Nomen gemeint (z.\,B. \textit{etwas Neues})? Dann handelt es sich um Substantivierung.
\item Bestimme für die Ellipse: \textbf{Kasus, Genus, Numerus, Deklinationstyp}; wähle danach die Endung.
\item Ist die Referenz unklar oder klingt der Satz unnatürlich? Wiederhole das Nomen.
\end{enumerate}

\section{Übungen mit Lösungen}
\textbf{Ergänze die richtige Form.}
\begin{enumerate}
\item Zwei Kleider: Das rot\_\_\_ gefällt mir. 
\item Zwei Mäntel: Ich kaufe den schwarz\_\_\_.
\item Zwei Taschen: Ich gehe mit der klein\_\_\_.
\item Zwei Computer: Mein alt\_\_\_ ist langsamer als mein neu\_\_\_.
\item Zwei Fahrräder: Welches gehört dir? -- Das blau\_\_\_.
\end{enumerate}
\textbf{Lösungen:} 1. rote; 2. schwarzen; 3. kleinen; 4. alter, neuer; 5. blaue.

\section{Verwandte Themen}
Adjektivdeklination; Nominalellipse; Nominalisierung/Substantivierung von Adjektiven; Pronomen und pronominale Artikelwörter; Kasus; Groß- und Kleinschreibung.
\end{document}
